% sec-06-cmc.tex - Chemistry, Manufacturing and Control (final stage)
\section{Chemistry, Manufacturing and Control Information}

\subsection{Chemistry, Manufacturing and Control}

This section presents the Chemistry, Manufacturing and Control (CMC)
information for the investigational drug daraxonrasib (RMC-6236) supporting the
Phase 1, first-in-human, open-label, single-arm investigation described
elsewhere in this Investigational New Drug application (IND v1.0, repository
v4.3.0). It is organized and submitted in accordance with 21 CFR \S
312.23(a)(7), which requires sufficient information about the composition,
manufacture, and control of the drug substance and the drug product to assure
the proper identification, quality, purity, and strength of the investigational
drug, scaled to the early phase of investigation, the limited number of
participants, and the short duration of exposure characteristic of a Phase 1
study \cite{phase1guidance}. Consistent with the FDA guidance on the content
and format of CMC information for a Phase 1 study, the emphasis here is on the
information needed to assure participant safety, including identity, strength,
quality, purity, and stability, rather than on the detailed validation data
appropriate to later phases \cite{phase1guidance}. 

The graded approach embodied
in 21 CFR \S 312.23(a)(7)(i) recognizes that the amount of CMC information
needed varies with the phase of investigation, the proposed duration, the dosage
form, and the amount of information otherwise available; this submission supplies
the identity, strength, quality, purity, and stability data appropriate to a
short perioperative dosing schedule in up to eighteen treated participants at a
single academic site, with approximately thirty-six participants screened to
reach that treated population \cite{phase1guidance}. As the investigation
proceeds and the manufacturing process and analytical methods mature, the CMC
information will be updated through information amendments under 21 CFR \S
312.31 as required, and any change that materially affects the safety of the
investigational product, alters the intended strengths, or modifies the
container-closure system will be reported before the affected product is
administered to participants.

The investigational regimen couples this oral drug with an on-premises large
language model (LLM) directed eight-arm robotic pancreaticoduodenectomy
(Whipple) platform that is the subject of the parallel investigational device
exemption content of this submission. The device is a durable instrument and is
not a drug, a biologic, or a consumable manufactured lot; its design,
specification, calibration, and quality controls are addressed in the device
description and the Physical AI System Description fields of 21 CFR \S
312.23(g), not in this CMC section
\cite{daraxwhipple,onpremwhippl}. The clean separation of the durable robotic
instrument and its hash-pinned advisory model from the manufactured drug
substance and drug product is deliberate: the verification-before-generation
discipline that governs the device, including its ten-gate ACCEPT, BLOCK, and
ESCALATE suite and its alignment with the Verification Before Generation in
Physical AI Oncology Trials Act of 2026, operates on software and kinematics and
has no bearing on the chemical identity or pharmaceutical quality of the tablet
\cite{congress9510,clintrustpai}. Accordingly, the four CMC subsections below
address only the orally administered drug substance and drug product, the
non-applicable placebo, and labeling, followed by the environmental assessment.
The CMC information architecture and the mapping of each subsection to its
content is shown in Figure 17.

\needspace{8\baselineskip}
\begin{mermaidfig}
\node[mmgoal] (cmc) at (5.2,0)      {Chemistry, Manufacturing\\and Control ($\S$7)};
\node[mmproc] (ds)  at (0,-2.8)     {7.1.1 Drug substance\\daraxonrasib\\(RMC-6236)};
\node[mmproc] (dp)  at (4.2,-2.8)   {7.1.2 Drug product\\oral tablet,\\160 / 220 / 300 mg};
\node[mmctx]  (pl)  at (8.4,-2.8)   {7.1.3 Placebo product\\not applicable\\(open-label)};
\node[mmproc] (lb)  at (12.6,-2.8)  {7.1.4 Labeling\\``Caution:\\investigational''\\21 CFR 312.6};
\node[mmproc] (ea)  at (12.6,0)     {7.2 Environmental\\assessment /\\categorical\\exclusion 25.31};
\node[mmin]   (ds1) at (0,-5.3)     {Characterization,\\specification,\\impurities};
\node[mmin]   (dp1) at (4.2,-5.3)   {Composition,\\container-closure,\\child-resistant};
\node[mmin]   (st)  at (8.4,-5.3)   {Stability program\\(supports shelf life)};
\draw[mmedge] (cmc) -- (ds);
\draw[mmedge] (cmc) -- (dp);
\draw[mmedge] (cmc) -- (pl);
\draw[mmedge] (cmc) -- (lb);
\draw[mmedge] (cmc.east) to (ea.west);
\draw[mmedge] (ds) -- (ds1);
\draw[mmedge] (dp) -- (dp1);
\draw[mmedge] (dp.south east) to (st.north);
\end{mermaidfig}
\vspace{-0.65cm}
\figcaption{Figure 17. The drug substance (daraxonrasib, RMC-6236) is described
by characterization, specification, and impurity controls; the drug product is
the oral tablet at the 160, 220, and 300 mg strengths with its composition\\and
child-resistant, light-protective container-closure under 21 CFR $\S$312.6,
supported by a stability program;\\the placebo product is not applicable for this
open-label study; labeling carries the investigational caution\\statement; and an
environmental assessment or categorical exclusion under 21 CFR $\S$25.31 is
provided.}

The remainder of \S 7.1 follows the drug substance through the drug product to
labeling, with each manufactured attribute tied to the test, analytical
procedure, and acceptance criterion that controls it. Throughout, the governing
principle is that every release and stability acceptance criterion exists to
assure one of the five Phase 1 quality attributes named in 21 CFR \S
312.23(a)(7)(i), namely identity, strength (potency), quality, purity, and
stability, and that any attribute not yet supported by a finalized validated
method is carried as a provisional interim Phase 1 specification that will be
tightened, not loosened, as data accumulate \cite{phase1guidance}. Table 7.1
maps each of these five attributes to the drug-substance and drug-product
controls described below, so that a reviewer can trace any single quality concern
from the regulatory requirement to the specific test that addresses it.

\needspace{10\baselineskip}
\begin{table}[t]
\centering
\caption{Mapping of the five Phase 1 CMC quality attributes of 21 CFR \S
312.23(a)(7)(i) to the drug-substance and drug-product controls described in this
section.}
\label{tab:cmc-attribute-map}
\begin{tabularx}{\textwidth}{L{3.0cm} L{5.5cm} Y}
\toprule
\textbf{Quality attribute} & \textbf{Drug-substance control} & \textbf{Drug-product control} \\
\midrule
Identity & IR, UV to visible, HPLC retention time, NMR, HRMS, chiral HPLC against reference standard & HPLC retention time plus a spectroscopic identity test on the finished tablet \\
Strength (potency) & HPLC assay, 98.0 to 102.0 percent on anhydrous, solvent-free basis & HPLC assay, 95.0 to 105.0 percent of label claim; content uniformity by USP \\
Quality & Solid-state form, particle size, water, residue on ignition, microbial limits & Appearance, hardness, friability, disintegration, dissolution, microbial limits \\
Purity & Organic impurities (ICH Q3A), genotoxic impurities (ICH M7), residual solvents (ICH Q3C), elemental impurities (ICH Q3D) & Degradation products (ICH Q3B) controlled to interim Phase 1 limits \\
Stability & Parallel ICH Q1A program establishing the retest period & ICH Q1A and Q1B program supporting the assigned shelf life \\
\bottomrule
\end{tabularx}
\vspace*{-0.2cm}
\figcaption{Table 7.1. Each regulatory quality attribute is addressed by at least
one drug-substance control and one drug-product control. Interim Phase 1
acceptance criteria are scaled to the early phase of investigation.}
\end{table}

Two cross-cutting commitments frame the CMC information that follows. First, every
acceptance criterion presented in this section is an interim Phase 1 criterion: it
is set to assure participant safety at the current state of process and analytical
knowledge and is expected to tighten, never to loosen, as additional batch data,
finalized validated methods, and longer-term stability data accumulate. Second,
any change to the drug substance or drug product that could affect the safety,
identity, strength, quality, purity, or stability of the investigational product
is managed by formal change control and reported to the IND by information
amendment under 21 CFR \S 312.31 before the affected material is administered to
participants. Table 7.11, at the close of \S 7.1, lists the principal CMC changes
that trigger an information amendment together with the regulatory basis and the
timing of the report; the remainder of this subsection should be read against
that change-control framework, so that each specification and each control is
understood as the current, amendable expression of an evolving but always
safety-protective body of CMC knowledge \cite{phase1guidance}.

\subsubsection{Drug Substance}

\paragraph{Nomenclature, description, and pharmacological class.}
The drug substance is daraxonrasib, also designated by the developmental code
RMC-6236, an orally bioavailable small molecule of the RAS(ON)
multi-selective (pan-RAS) inhibitor class. Daraxonrasib acts by binding the
active, guanosine triphosphate (GTP) bound state of mutant RAS in a tri-complex
with the endogenous chaperone cyclophilin A, thereby occluding the
RAS to effector interface and suppressing downstream mitogen-activated protein
kinase and phosphoinositide 3-kinase signaling across the common KRAS
oncoprotein variants found in pancreatic ductal adenocarcinoma (PDAC), including
the G12D, G12V, and G12R substitutions targeted by this protocol
\cite{daraxwhipple,pdac060s2030}. The agent received FDA Breakthrough Therapy
designation in June 2025 for previously treated metastatic PDAC bearing a KRAS
G12 mutation, and 300 mg administered orally once daily is the recommended
Phase 2 dose established in the pan-KRAS RASolute 302 program to which the dose
and exposure assumptions of this protocol are anchored
\cite{daraxwhipple,pdac060s2030}. The drug substance is a non-biologic, chirally
defined organic small molecule; it is not a radioactive drug, is not a
controlled substance, and is not a peptide, protein, or nucleic acid. The drug
substance is presented as a solid, and in the finished product it is formulated
as an immediate-release oral tablet. Table 7.2 collects the nomenclature,
classification, and regulatory descriptors of the drug substance in the form
expected for the drug-substance description of 21 CFR \S 312.23(a)(7).

\needspace{10\baselineskip}
\begin{table}[t]
\centering
\caption{Drug-substance nomenclature, classification, and regulatory descriptors
for daraxonrasib (RMC-6236).}
\label{tab:cmc-ds-nomenclature}
\begin{tabularx}{\textwidth}{L{5.1cm} Y}
\toprule
\textbf{Descriptor} & \textbf{Value} \\
\midrule
Nonproprietary name & Daraxonrasib \\
Developmental code & RMC-6236 \\
Pharmacological class & RAS(ON) multi-selective (pan-RAS) small-molecule inhibitor \\
Mechanism of action & Binds GTP-bound active state of mutant RAS in a tri complex with cyclophilin A; occludes RAS to effector interface \\
Targeted variants in this study & KRAS G12D, G12V, and G12R \\
Molecular nature & Non-biologic, chirally defined organic small molecule \\
Physical state & Solid (crystalline) \\
Route of administration & Oral, once daily \\
Controlled-substance status & Not a controlled substance \\
Radioactive-drug status & Not a radioactive drug \\
Regulatory milestone & FDA Breakthrough Therapy designation (June 2025), KRAS G12 metastatic PDAC \\
Recommended Phase 2 dose & 300 mg once daily (RASolute 302) \\
\bottomrule
\end{tabularx}
\vspace*{-0.2cm}
\figcaption{Table 7.2. Nomenclature and regulatory descriptors of the drug
substance. The substance is a\\small-molecule pan-RAS inhibitor and carries no
controlled-substance or radioactive-drug\\obligations; abuse-potential and
radioactive-drug considerations are addressed further in $\S$10.}
\end{table}

\paragraph{Manufacturer and method of preparation.}
The drug substance is manufactured by, or under contract to, the sponsor's
designated active pharmaceutical ingredient manufacturer under current good
manufacturing practice (cGMP) appropriate to the investigational stage, and is
released against the drug-substance specification summarized in Table 7.3. The
synthetic route is a defined, multi-step convergent organic synthesis with
isolated, characterized intermediates and controlled reagents, catalysts, and
solvents; in-process controls govern the critical steps, and the final isolation
and drying steps establish the solid-state form, particle size distribution, and
residual-solvent profile of the substance. The name, address, and
establishment information of the drug-substance manufacturer, the detailed
reaction scheme, the list of starting materials and reagents, the in-process
controls, and the batch records are provided in the supporting CMC documentation
on file with the sponsor and are cross-referenced here consistent with the
Phase 1 level of CMC detail \cite{phase1guidance}. 

A statement that the drug
substance is manufactured in compliance with cGMP, and the certificate of
analysis for the clinical batch, accompany each shipment. Consistent with the
Phase 1 graded approach, the submission identifies the regulatory starting
materials, describes the controls applied at each critical processing step, and
commits to reporting any change to the route of synthesis, the source of a
starting material, or the established solid-state form by information amendment
under 21 CFR \S 312.31 before such material is used to manufacture clinical drug
product \cite{phase1guidance}. Because solid-state form and particle size
distribution are determined at the final isolation and drying steps and in turn
govern dissolution and oral exposure, these two attributes are treated as
critical quality attributes of the drug substance and are controlled both as
in-process parameters and as release tests. 

The regulatory starting materials are
defined points in the synthesis at which the supplier, the specification, and the
control strategy are fixed, and the steps downstream of those starting materials
are conducted under cGMP appropriate to the investigational stage with isolated,
characterized intermediates; this defined boundary allows the impurity-control
strategy of ICH Q3A and the genotoxic-impurity assessment of ICH M7 to be applied
consistently from a known, controlled origin to the final drug substance
\cite{ichpaistduni}. The certificate of analysis that accompanies each clinical
batch reports the result of every test in the drug-substance specification against
its acceptance criterion, and the batch is released only when every result
conforms; non-conforming material is not released for use in drug
manufacture.

\paragraph{Characterization.}
The structure and identity of daraxonrasib are established by a panel of
orthogonal physicochemical methods. Structural elucidation and confirmation use
proton and carbon nuclear magnetic resonance spectroscopy, high-resolution mass
spectrometry, infrared spectroscopy, and ultraviolet to visible spectroscopy;
elemental analysis confirms empirical composition. Chiral purity is confirmed by
chiral high-performance liquid chromatography (HPLC) to assure the correct
stereochemical configuration of the defined stereocenters. The solid-state form
(polymorph) is characterized by X-ray powder diffraction and differential
scanning calorimetry, and the manufacturing process is controlled to deliver the
intended crystalline form, because solid-state form governs dissolution and
hence the oral exposure on which the pharmacokinetically gated restart advisory
of the protocol depends. Hygroscopicity, aqueous and physiological-pH
solubility, the partition coefficient, and the pKa are determined to support the
biopharmaceutic understanding of the tablet. 

These characterization data
constitute the reference standard against which clinical batches are compared.
The biopharmaceutic relevance of these properties is concrete in this protocol:
the perioperative pause-and-restart advisory compares the measured serum
daraxonrasib trough against a 0.5 ng/mL threshold, and across the 32-iteration
deterministic sweep the modeled baseline trough is 0.45 ng/mL, sub-threshold in
every iteration \cite{daraxwhipple}. Because that comparison turns on the oral
exposure delivered by the tablet, any drift in solid-state form, particle size,
or dissolution that lowered bioavailability would shift the modeled trough and
the advisory it informs; controlling these attributes at release therefore
protects not only pharmaceutical quality but the integrity of a safety-relevant
clinical decision. The advisory itself remains a second-opinion, human-confirmed
recommendation and never an autonomous dosing action, consistent with the
prohibition on full autonomy under 21 CFR \S 312.21(e) \cite{onpremwhippl}.

\paragraph{Analytical procedures and reference standard.}
The analytical procedures used to release and to test the stability of the drug
substance are the chromatographic, spectroscopic, and physical methods listed in
the specification of Table 7.3, applied at the Phase 1 level of method
qualification appropriate to a first-in-human study. The assay and
related-substances methods are reversed-phase HPLC procedures with ultraviolet
detection that have been demonstrated to be stability-indicating, that is, able to
resolve the parent drug substance from its known and potential degradation
products so that any increase in a degradant on storage is detected and
quantified; the chiral HPLC method resolves the defined stereoisomer from its
undesired enantiomer; and the headspace gas chromatography, inductively coupled
plasma mass spectrometry, Karl Fischer, and X-ray powder diffraction procedures
address residual solvents, elemental impurities, water content, and solid-state
form respectively. 

At the Phase 1 stage these procedures are qualified rather than
fully validated, with specificity, linearity over the working range, accuracy,
precision, and the limits of detection and quantitation established to the extent
needed to support the acceptance criteria; full method validation under ICH Q2 is
completed as the program advances and is reported by information amendment. A
qualified primary reference standard of daraxonrasib, fully characterized by the
orthogonal panel described above and assigned a potency on the anhydrous,
solvent-free basis, is established and is used to calibrate the assay and identity
tests; working reference standards are qualified against the primary standard.
Each clinical batch is compared against this reference standard, and the
certificate of analysis for each batch reports the result of every test in the
specification against its acceptance criterion. Where a new clinical batch is
introduced during the investigation, comparability against earlier batches is
assessed against the same specification and reference standard so that batch-to-
batch consistency is demonstrated before the new batch is administered.

\paragraph{Specification and control of impurities.}
The drug-substance specification, with tests, analytical procedures, and
acceptance criteria appropriate to a first-in-human Phase 1 study, is given in
Table 7.3. Organic impurities (process-related and degradation-related), residual
solvents, and elemental impurities are controlled to limits consistent with the
applicable International Council for Harmonisation (ICH) quality guidelines (ICH
Q3A for impurities in new drug substances, ICH Q3C for residual solvents, and
ICH Q3D for elemental impurities), as adapted to the limited duration and
participant exposure of a Phase 1 investigation
\cite{ichpaistduni,phase1guidance}. Any individual or total impurity exceeding
its qualification threshold is qualified by reference to the nonclinical
toxicology program summarized in the pharmacology and toxicology section of this
IND. Genotoxic-impurity risk is assessed and controlled under the ICH M7
framework, with structural alerts evaluated and relevant potential mutagenic
impurities controlled to the threshold of toxicological concern. The acceptance
criteria below are interim Phase 1 specifications and will be tightened as
process knowledge and additional batch data accumulate.

A worked example illustrates how the impurity acceptance criteria translate into
an absolute daily exposure at the doses studied. At the highest dose level, dose
level 3 (DL3), the once-daily dose is 300 mg; an individual organic impurity
controlled to not more than 0.5 percent therefore corresponds to at most 1.5 mg
of that impurity per day, and the not-more-than 2.0 percent total-impurity limit
corresponds to at most 6.0 mg of total impurities per day. At the starting dose
level 1 (DL1) dose of 160 mg the same percentage limits correspond to at most
0.8 mg per individual impurity and 3.2 mg total per day, and at dose level 2
(DL2) the 220 mg dose corresponds to at most 1.1 mg and 4.4 mg respectively.
These absolute daily amounts, taken across the short perioperative dosing
schedule, are the quantities that the nonclinical qualification is required to
support, and they remain within the qualification framework of ICH Q3A as
adapted for a Phase 1 exposure of limited duration in a population of up to
eighteen treated participants \cite{ichpaistduni,phase1guidance}. Table 7.4
records these worked daily impurity exposures at each dose level.

\needspace{12\baselineskip}
\begin{table}[t]
\centering
\caption{Drug-substance specification for daraxonrasib (RMC-6236): interim
Phase 1 release and stability tests, analytical procedures, and acceptance
criteria.}
\label{tab:cmc-ds-spec}
\begin{tabularx}{\textwidth}{L{4.1cm} L{4.0cm} Y}
\toprule
\textbf{Attribute / test} & \textbf{Analytical procedure} & \textbf{Acceptance criterion (interim Phase 1)} \\
\midrule
Appearance & Visual & White to off-white solid; free of foreign matter \\
Identity (spectroscopic) & IR; UV to visible & Conforms to reference standard \\
Identity (chromatographic) & HPLC retention time & Retention time matches reference standard \\
Identity (structural) & 1H and 13C NMR; HRMS & Consistent with assigned structure \\
Stereochemical purity & Chiral HPLC & Single defined stereoisomer; undesired enantiomer not more than 0.5 percent \\
Assay (content) & HPLC with UV detection & 98.0 to 102.0 percent on the anhydrous, solvent-free basis \\
Organic impurities (each) & HPLC, related substances & Not more than 0.5 percent per specified or unspecified impurity (ICH Q3A) \\
Organic impurities (total) & HPLC, related substances & Not more than 2.0 percent total \\
Genotoxic impurities & LC-MS or GC-MS, ICH M7 & Below the threshold of toxicological concern for specified mutagenic impurities \\
Residual solvents & Headspace GC & Within ICH Q3C Class limits for solvents used in synthesis \\
Elemental impurities & ICP-MS & Within ICH Q3D oral permitted daily exposure limits \\
Water content & Karl Fischer titration & Not more than 1.0 percent w/w \\
Residue on ignition / sulfated ash & USP general chapter & Not more than 0.2 percent \\
Solid-state form & X-ray powder diffraction & Conforms to the designated reference polymorph \\
Particle size distribution & Laser diffraction & Within the validated process range supporting dissolution \\
Microbial limits & USP total aerobic / yeast and mold & Within USP limits for an oral non-sterile substance \\
\bottomrule
\end{tabularx}
\vspace*{-0.2cm}
\figcaption{Table 7.3. Interim Phase 1 drug-substance specification. Acceptance
criteria are scaled to\\the early phase of investigation and will be tightened as
batch and process data accumulate;\\impurity limits reference ICH Q3A, Q3C, Q3D,
and M7 as adapted for Phase 1.}
\end{table}

\needspace{9\baselineskip}
\begin{table}[t]
\centering
\caption{Worked daily impurity exposure at each 3+3 dose level, derived from the
interim percentage acceptance criteria of the drug-substance specification.}
\label{tab:cmc-impurity-exposure}
\begin{tabularx}{\textwidth}{L{2.6cm} C{2.2cm} C{3.4cm} Y}
\toprule
\textbf{Dose level} & \textbf{Daily dose} & \textbf{Max single impurity (0.5 percent)} & \textbf{Max total impurities (2.0 percent)} \\
\midrule
DL1 (start) & 160 mg & 0.8 mg per day & 3.2 mg per day \\
DL2 & 220 mg & 1.1 mg per day & 4.4 mg per day \\
DL3 (RP2D anchor) & 300 mg & 1.5 mg per day & 6.0 mg per day \\
\bottomrule
\end{tabularx}
\vspace*{-0.3cm}
\figcaption{Table 7.4. Worked impurity exposures. The percentage acceptance
criteria of Table 7.3 translate\\into the absolute daily amounts shown; these are
the quantities the nonclinical qualification\\under ICH Q3A is required to support
across the short perioperative dosing schedule.}
\end{table}

\subsubsection{Drug Product}

\paragraph{Description and composition.}
The drug product is an immediate-release, film-coated oral tablet supplied in the
three strengths required to compose the once-daily doses of the 3+3 escalation:
160 mg for dose level 1 (DL1), 220 mg for dose level 2 (DL2), and 300 mg for
dose level 3 (DL3), the last of which corresponds to the RASolute 302 recommended
Phase 2 dose \cite{daraxwhipple,pdac060s2030}. The three strengths are
dose-proportional, sharing a common qualitative composition and a common blend so
that the only quantitative difference among strengths is the content of drug
substance and the corresponding fill weight, with the tablets differentiated by
size, debossing, or both to prevent dose confusion in this open-label study. The
qualitative and representative quantitative composition is given in Table 7.5. The
excipients are standard, compendial pharmaceutical excipients used at
conventional levels; the formulation contains a filler or diluent, a binder, a
disintegrant, a glidant, a lubricant, and a non-functional aqueous film coat. No
excipient of human or animal origin presenting a transmissible spongiform
encephalopathy risk is used, and no novel excipient is included. Because no novel
excipient is present, no separate excipient safety qualification is required, and
the excipient information is provided at the Phase 1 level of detail
\cite{phase1guidance}. The functional roles of the excipients follow conventional
immediate-release design: the microcrystalline cellulose and the second filler
provide bulk and compactibility to the target fill weight; the binder confers
granule and tablet cohesion; the croscarmellose sodium disintegrant promotes rapid
break-up in the gastrointestinal tract to support the provisional not-less-than 80
percent dissolved in 30 minutes immediate-release criterion; the colloidal silicon
dioxide glidant assures uniform powder flow into the press, which protects weight
uniformity and hence content uniformity; the magnesium stearate lubricant, held to
a conventional low level, prevents sticking and picking without retarding
disintegration; and the hypromellose-based aqueous film coat is cosmetic and
non-functional, applied to a target weight gain and not intended to modify release.
Because the coat is non-functional, its only release-relevant requirement is that
it not retard dissolution, which is confirmed by the dissolution test of the
finished-product specification.

The single-strength, dose-proportional, common-blend design has a direct safety
and accountability benefit in an open-label 3+3 escalation. Because the three
strengths differ only in drug-substance content and fill weight, a single
validated blend and a single set of in-process controls support all three
presentations, and the physical differentiation by size and debossing provides a
visible, dispensing-pharmacy-verifiable distinction among the assigned dose
levels so that a participant at DL1 cannot be inadvertently dispensed a DL3
tablet. Table 7.6 summarizes the three strengths, their dose levels, and the
identifying features that distinguish them at the investigational-product
pharmacy.

\needspace{12\baselineskip}
\begin{table}[t]
\centering
\caption{Representative qualitative and quantitative composition of the
daraxonrasib immediate-release oral tablet (per the 160 mg strength; the 220 mg
and 300 mg strengths are dose-proportional from a common blend).}
\label{tab:cmc-dp-comp}
\begin{tabularx}{\textwidth}{L{4.4cm} L{4.6cm} Y}
\toprule
\textbf{Component} & \textbf{Function} & \textbf{Quantity (160 mg strength) and notes} \\
\midrule
Daraxonrasib (RMC-6236) & Active pharmaceutical ingredient & 160 mg per tablet (220 mg / 300 mg in higher strengths) \\
Microcrystalline cellulose & Filler / diluent & Quantity sufficient to target fill weight \\
Lactose monohydrate or mannitol & Filler / diluent & Compendial grade, conventional level \\
Povidone or hypromellose & Binder & Conventional level \\
Croscarmellose sodium & Disintegrant & Conventional level for immediate release \\
Colloidal silicon dioxide & Glidant & Conventional level \\
Magnesium stearate & Lubricant & Conventional level (typically not more than 1.0 percent) \\
Film coat (hypromellose-based) & Aqueous non-functional coat & Applied to target weight gain; removed as water on drying \\
Purified water & Coating vehicle & Removed during processing (not present in finished tablet) \\
\bottomrule
\end{tabularx}
\vspace*{-0.25cm}
\figcaption{Table 7.5. Drug-product composition. All excipients are compendial and
used at conventional levels; the film coat\\is non-functional and cosmetic. The
220 mg and 300 mg strengths are manufactured dose-proportionally from a common
blend, differing in drug-substance content and fill weight and in tablet
size or debossing for identification.}
\end{table}

\needspace{9\baselineskip}
\begin{table}[t]
\centering
\caption{Drug-product strengths mapped to the 3+3 dose levels, with the
identifying features that distinguish them at dispensing.}
\label{tab:cmc-dp-strengths}
\begin{tabularx}{\textwidth}{L{1.2cm} C{2.1cm} L{2.7cm} Y}
\toprule
\textbf{Strength} & \textbf{Dose level} & \textbf{Once-daily dose} & \textbf{Identification feature} \\
\midrule
160 mg & DL1 (start) & 160 mg & Smallest tablet; distinct debossing; verified at dispense \\
220 mg & DL2 & 220 mg & Intermediate tablet size; distinct debossing \\
300 mg & DL3 & 300 mg & Largest tablet; distinct debossing; equals RASolute 302 recommended Phase 2 dose \\
\bottomrule
\end{tabularx}
\vspace*{-0.2cm}
\figcaption{Table 7.6. Strength-to-dose-level mapping. The dose-proportional
common-blend design and\\the physical differentiation among strengths support
accurate dispensing and drug accountability\\in this open-label study; no
strength carries any masking obligation because the study is open-label.}
\end{table}

\paragraph{Manufacture.}
The drug product is manufactured under cGMP appropriate to the investigational
stage by the sponsor's designated drug-product manufacturer. The process is a
conventional solid-oral-dosage process: dispensing and screening of the drug
substance and excipients, blending (with wet or dry granulation as defined in the
batch record), lubrication, compression of the final blend into core tablets on a
rotary tablet press under in-process controls for tablet weight, hardness,
thickness, and friability, and aqueous film coating of the cores in a perforated
coating pan to a target weight gain, followed by drying. In-process controls at
compression include tablet weight and weight uniformity, hardness, thickness,
disintegration, and visual inspection; the coating step is controlled for weight
gain and appearance. The batch formula, the master and executed batch records,
the in-process control limits, and the name and address of the drug-product
manufacturer are provided in the supporting CMC documentation and are
cross-referenced here at the Phase 1 level of detail \cite{phase1guidance}. A
certificate of analysis accompanies each clinical batch, and each batch is
released against the drug-product specification before use. Table 7.7 lists the
principal in-process controls applied during manufacture, the manufacturing step
at which each is exercised, and its purpose, so that a reviewer can see how
finished-product quality is built in rather than tested in.

\needspace{11\baselineskip}
\begin{table}[t]
\centering
\caption{Principal in-process controls applied during manufacture of the
daraxonrasib oral tablet.}
\label{tab:cmc-ipc}
\begin{tabularx}{\textwidth}{L{3.4cm} L{3.6cm} Y}
\toprule
\textbf{In-process control} & \textbf{Manufacturing step} & \textbf{Purpose} \\
\midrule
Blend uniformity & Blending / lubrication & Assures uniform drug-substance distribution before compression, supporting content uniformity \\
Tablet weight and weight uniformity & Compression & Controls dose delivered per tablet; supports assay and uniformity of dosage units \\
Hardness & Compression & Assures mechanical integrity and consistent disintegration and dissolution \\
Thickness & Compression & Confirms consistent tooling and compaction \\
Friability & Compression & Limits chipping and dust during coating and packaging \\
Disintegration & Compression & Early surrogate confirming immediate-release behavior \\
Coating weight gain & Film coating & Controls cosmetic coat to target; non-functional coat must not retard release \\
Appearance & Coating / inspection & Confirms intact, defect-free, correctly identified tablets \\
\bottomrule
\end{tabularx}
\vspace*{-0.2cm}
\figcaption{Table 7.7. In-process controls. Quality is built into the product at
each step; the listed controls govern the critical compression and coating
operations and feed forward to the finished-product release tests of Table 7.8.}
\end{table}

\paragraph{Drug-product specification.}
Each clinical batch is tested and released against a drug-product specification
that includes appearance, identity of the drug substance (by HPLC retention time
and a spectroscopic method), assay (95.0 to 105.0 percent of label claim by HPLC),
content uniformity by the USP uniformity-of-dosage-units test, degradation
products controlled to interim Phase 1 limits consistent with ICH Q3B (not more
than 0.5 percent per specified or unspecified degradant and not more than 2.0
percent total), dissolution by the USP apparatus and medium validated for the
product (for example, not less than 80 percent dissolved in 30 minutes as a
provisional immediate-release criterion to be confirmed by the developing
dissolution method), water content, and microbial limits within USP acceptance
criteria for an oral non-sterile product. The dissolution method and acceptance
criterion are provisional and will be finalized as biopharmaceutic and stability
data accrue. The full interim Phase 1 drug-product release and stability
specification is collected in Table 7.8, with each attribute paired to its
analytical procedure and its acceptance criterion; the same panel of
stability-indicating attributes is monitored on the stability program described
below, so that any change on storage is detected against the same criteria that
govern release.

\needspace{12\baselineskip}
\begin{table}[t]
\centering
\caption{Drug-product specification for the daraxonrasib oral tablet: interim
Phase 1 release and stability tests, analytical procedures, and acceptance
criteria.}
\label{tab:cmc-dp-spec}
\begin{tabularx}{\textwidth}{L{4.0cm} L{4.0cm} Y}
\toprule
\textbf{Attribute / test} & \textbf{Analytical procedure} & \textbf{Acceptance criterion (interim Phase 1)} \\
\midrule
Appearance & Visual & Film-coated tablet of defined color and shape; correct debossing; defect-free \\
Identity (chromatographic) & HPLC retention time & Matches reference standard \\
Identity (spectroscopic) & UV to visible or IR & Conforms to reference standard \\
Assay (label claim) & HPLC with UV detection & 95.0 to 105.0 percent of label claim \\
Content uniformity & USP uniformity of dosage units & Meets USP acceptance value \\
Degradation products (each) & HPLC, stability-indicating & Not more than 0.5 percent per specified or unspecified degradant (ICH Q3B) \\
Degradation products (total) & HPLC, stability-indicating & Not more than 2.0 percent total \\
Dissolution & USP apparatus, validated medium & Not less than 80 percent dissolved in 30 minutes (provisional) \\
Water content & Karl Fischer titration & Within the limit supporting stability \\
Microbial limits & USP total aerobic / yeast and mold & Within USP limits for an oral non-sterile product \\
\bottomrule
\end{tabularx}
\vspace*{-0.2cm}
\figcaption{Table 7.8. The dissolution
method and limit are provisional pending the developing method; degradation
limits reference ICH Q3B; the same stability-indicating attributes are monitored
on the stability program of Table 7.10.}
\end{table}

\paragraph{Container-closure system.}
The drug product is packaged in a child-resistant, light-protective,
moisture-protective container-closure system, namely high-density polyethylene
bottles with a child-resistant, induction-sealed closure and an appropriate
desiccant, or equivalent child-resistant unit-dose blister packaging. The
container-closure components contacting the product are of compendial
pharmaceutical grade and are suitable for an oral solid dosage form; the
child-resistant feature satisfies the special-packaging expectation for an
investigational oral oncology agent, and the light-protective and
moisture-protective features support the stability claim. Packaging and labeling
operations are performed under cGMP, and the investigational caution statement and
the label content described in Table 7.9 and in $\S$7.1.4 appear on the immediate
and any secondary container in accordance with 21 CFR \S 312.6 \cite{phase1guidance}.
Each container-closure feature is tied to the quality attribute it protects: the
high-density polyethylene bottle and induction seal together with the desiccant
control moisture ingress and so protect assay and degradation-product limits on
storage; the light-protective feature, qualified by the ICH Q1B photostability
challenge, protects against photodegradation; and the child-resistant closure
satisfies the special-packaging expectation for an oral oncology agent dispensed
for at-home once-daily administration during the perioperative schedule.

\paragraph{Stability program.}
Stability of the drug product is supported by an ongoing stability program
conducted under ICH Q1A conditions on representative batches in the proposed
container-closure system, using stability-indicating analytical methods. The
program comprises long-term, intermediate, and accelerated conditions, with
attributes monitored at each station including appearance, assay, degradation
products, dissolution, and water content; a photostability evaluation under ICH
Q1B and an in-use stability evaluation for the multi-dose bottle presentation are
included. The program and its current status are summarized in Table 7.10. The
provisional shelf life (expiry or retest period) assigned to the clinical drug
product is supported by the available long-term and accelerated data and is
applied through expiry dating on the label; product is dispensed only within the
labeled expiry or retest date. Stability is monitored throughout the
investigation, and the shelf life will be extended through information amendments
as additional long-term data become available \cite{phase1guidance}. The drug
substance is also placed on a parallel ICH Q1A stability program to establish its
retest period. 

The accelerated arm at 40 degrees Celsius and 75 percent relative
humidity provides early detection of degradation pathways and supports provisional
dating and the assessment of any temperature excursion incurred during shipment or
site storage; should a significant change occur at the accelerated condition, the
intermediate arm at 30 degrees Celsius and 65 percent relative humidity is read to
characterize the behavior under the intermediate condition, consistent with the
ICH Q1A decision framework. The in-use stability evaluation is specifically
relevant to the multi-dose bottle presentation dispensed for at-home once-daily
administration during the perioperative schedule: it establishes the period after
first opening during which the desiccated, light-protective bottle continues to
meet the assay, degradation-product, dissolution, and water-content criteria, and
that labeled in-use period is reflected in the storage and handling directions of
$\S$7.1.4. 

Because the perioperative regimen pauses dosing before surgery and
restarts it only on the advisory-determined day, a participant's bottle may be
opened, set aside during the operative and immediate postoperative interval, and
returned to use at restart; the in-use program is designed to cover precisely this
open-then-resume pattern, and the unit-dose blister alternative is available where
an even more conservative moisture control is preferred. Any temperature excursion
incurred during shipment or site storage is assessed against the accelerated-arm
data under a documented excursion protocol, and product is returned to use only
when the excursion is shown to fall within the supported range; otherwise the
affected units are quarantined and destroyed under the drug-accountability
procedure.

\needspace{12\baselineskip}
\begin{table}[t]
\centering
\caption{Drug-product stability program for the daraxonrasib oral tablet:
storage conditions, time points, attributes monitored, and purpose.}
\label{tab:cmc-stability}
\begin{tabularx}{\textwidth}{L{2.3cm} L{3.5cm} L{3.9cm} Y}
\toprule
\textbf{Condition} & \textbf{Storage} & \textbf{Time points (months)} & \textbf{Attributes monitored / purpose} \\
\midrule
Long-term (ICH Q1A) & 25 C / 60 percent RH & 0, 3, 6, 9, 12, 18, 24 & Appearance, assay, degradants, dissolution, water; supports assigned shelf life \\
Intermediate (ICH Q1A) & 30 C / 65 percent RH & 0, 6, 12 & Same panel; triggered if a significant change occurs at accelerated \\
Accelerated (ICH Q1A) & 40 C / 75 percent RH & 0, 1, 2, 3, 6 & Same panel; supports provisional dating and excursion assessment \\
Photostability (ICH Q1B) & ICH option 1 light exposure & Single challenge & Confirms need for light-protective packaging \\
In-use stability & 25 C / 60 percent RH, bottle opened & Per protocol after first opening & Supports the labeled in-use period for the multi-dose bottle \\
Temperature excursion & Per documented excursion & As incurred & Supports continued-use decisions for monitored excursions \\
\bottomrule
\end{tabularx}
\vspace*{-0.2cm}
\figcaption{Table 7.10. Stability-indicating
methods are used at every station; the provisional shelf life is supported by the
long-term and accelerated data and is extended by information amendment as
additional long-term data accrue.}
\end{table}

\subsubsection{Placebo Product}

This is an open-label, single-arm, first-in-human study; there is no comparator
arm, no masking of treatment assignment, and consequently no placebo product. The
absence of a placebo is consistent with the dose-finding and early-feasibility
objectives of a Phase 1 investigation and with ICH E9 principles for early-phase
trials \cite{ichpaistduni,phase1guidance}. In accordance with the ReGARDD IND
content map, this subsection is retained for completeness and to preserve the
table-of-contents numbering, and the placebo product is recorded here as not
applicable. The bias-control measures appropriate to an open-label design,
including the masking of the pathologist who determines R0 resection status and
the central blinded radiographic reader, are described in study protocol and
in $\S$6.1 of this IND rather than through a placebo. 

These outcome-assessor
masking measures are methodological controls on endpoint adjudication, not
pharmaceutical products, and accordingly impose no manufacturing,
container-closure, or labeling obligation under this CMC section; in particular,
no placebo-matching tablet is manufactured, and the drug-product labeling of
$\S$7.1.4 is constructed so that nothing appearing on labeling provided to a
masked assessor reveals the assigned dose level or any outcome. Should a future
amendment introduce a randomized or placebo-controlled component, a corresponding
placebo product description, composition, and matching strategy, including the
appearance, debossing, and container-closure matching needed to preserve any
masking, would be submitted by information amendment under 21 CFR \S 312.31.

\subsubsection{Labeling}

The investigational drug product is labeled in accordance with 21 CFR \S 312.6,
which provides that the immediate package of an investigational drug intended for
human use bear a label with the statement ``Caution: New Drug. Limited by Federal
(or United States) law to investigational use,'' and that the label and labeling
bear no false or misleading statement and make no claim that the drug is safe or
effective for the purposes under investigation \cite{phase1guidance}. Because the
study is open-label, the participant-facing label identifies the product and the
assigned once-daily dose level; no information that would unblind a masked outcome
assessor (for example, the masked pathologist or the central radiographic reader)
appears on any label or in any labeling provided to those assessors. 

The required
content of the immediate-container label and of the supplied labeling is
summarized in Table 7.9. The label is affixed during cGMP packaging operations and
is verified at the investigational-product pharmacy against the dispensing record
before any unit is released to a participant. The investigational device and its
single-use accessories carry the investigational-device caution statement under
21 CFR part 812 and are not relabeled or repackaged at the site; their labeling is
addressed in the device content of this submission. The directions for use on the
label reinforce the perioperative dosing schedule: the once-daily dose is paused
before surgery and restarted only on the day determined by the human-confirmed,
LLM-bound restart advisory described in the protocol, so the label instructs the
participant to follow the study team's dosing direction and not to resume dosing
on their own, and it carries the ``swallow whole; do not crush, split, or chew''
instruction appropriate to a film-coated immediate-release tablet.

\needspace{13\baselineskip}
\begin{table}[t]
\centering
\caption{Investigational-product label and labeling content for daraxonrasib
oral tablet under 21 CFR \S 312.6.}
\label{tab:cmc-labeling}
\begin{tabularx}{\textwidth}{L{4.6cm} Y}
\toprule
\textbf{Label element} & \textbf{Content and basis} \\
\midrule
Investigational caution statement & ``Caution: New Drug. Limited by Federal (or United States) law to investigational use'' (21 CFR \S 312.6) \\
Product identity and strength & Daraxonrasib (RMC-6236) oral tablet; 160 mg, 220 mg, or 300 mg as assigned \\
Protocol and IND reference & Protocol number and IND number; sponsor protocol identifier \\
Lot and batch identification & Manufacturing lot number and a unique container identifier for accountability \\
Storage conditions & ``Store at controlled room temperature; protect from light and moisture; keep in original container'' \\
Expiry or retest date & Date supported by the stability program; not dispensed beyond this date \\
Quantity and directions & Net contents and the once-daily directions for use; ``Swallow whole; do not crush, split, or chew'' \\
Child-resistant statement & ``Keep out of reach of children''; child-resistant closure per special packaging \\
Sponsor identification & Name and address of the sponsor-investigator (ChemicalQDevice) \\
Unblinding control & No information appearing on labeling provided to masked assessors that would reveal dose level or outcome \\
\bottomrule
\end{tabularx}
\vspace*{-0.2cm}
\figcaption{Table 7.9. Label and labeling content. The immediate package bears the
21 CFR \S 312.6 investigational\\caution statement and makes no claim of safety or
effectiveness; the label supports drug\\accountability and protects the masked
outcome assessors of this open-label study.}
\end{table}

To close \S 7.1, the change-control framework introduced above is set out
explicitly. The CMC information in this section reflects the manufacturing process
and analytical methods as currently established for the Phase 1 investigation, and
it will evolve as process knowledge accrues, as methods are fully validated under
ICH Q2, and as longer-term stability data become available. Every such evolution
is governed by formal change control, and any change that could affect the safety
of the investigational product, alter the strengths supplied, modify the
container-closure system, or change the established route of synthesis or
solid-state form is reported to the IND by information amendment under 21 CFR \S
312.31 before the affected material is administered to participants
\cite{phase1guidance}. Table 7.11 lists the principal CMC changes that trigger an
information amendment, the regulatory basis for each, and the timing of the
report. This change-control discipline mirrors, in the pharmaceutical domain, the
re-validation discipline that governs the device side of this combined submission,
where any change to the robotic platform or its hash-pinned advisory model
requires the Phase 0 simulation-validation gate to be re-cleared before further
patient-facing activity; in both domains a controlled change is permitted only
after the relevant safety-protective evidence has been re-established
\cite{onpremwhippl}.

\clearpage

\needspace{12\baselineskip}
\begin{table}[t]
\centering
\caption{Principal CMC changes that trigger an information amendment under 21 CFR
\S 312.31, with regulatory basis and reporting timing.}
\label{tab:cmc-change-control}
\begin{tabularx}{\textwidth}{L{4.4cm} L{3.6cm} Y}
\toprule
\textbf{CMC change} & \textbf{Regulatory basis} & \textbf{Reporting timing} \\
\midrule
Change to the route of synthesis or a starting-material source & 21 CFR \S 312.31; ICH Q3A & Reported by information amendment before affected drug substance is used \\
Change to the established solid-state form or particle size process range & 21 CFR \S 312.31 & Reported before affected material is administered; supported by dissolution data \\
Tightening of an interim acceptance criterion & 21 CFR \S 312.31 & Reported at the next information amendment as data support tightening \\
Introduction of a new clinical batch & 21 CFR \S 312.31 & Comparability assessed against specification and reference standard before administration \\
Change to the container-closure system & 21 CFR \S 312.31 & Reported before the new system is used; supported by stability bridging \\
Extension of the assigned shelf life & 21 CFR \S 312.31 & Reported as additional long-term stability data accrue \\
Finalization of a validated method (for example, dissolution) & 21 CFR \S 312.31; ICH Q2 & Reported when the validated method and its criterion replace the provisional one \\
Introduction of a placebo or comparator (future amendment) & 21 CFR \S 312.31 & Full placebo description, composition, and matching strategy submitted before use \\
\bottomrule
\end{tabularx}
\vspace*{-0.2cm}
\figcaption{Table 7.11. CMC change control. Each listed change is managed by
formal change control and reported\\by information amendment under 21 CFR
$\S$312.31 before the affected material is administered, so\\that the CMC
information remains current and safety-protective throughout the investigation.}
\end{table}

\subsection{Environmental Assessment}

The sponsor claims a categorical exclusion from the requirement to prepare an
environmental assessment or an environmental impact statement under 21 CFR \S
25.31(e), which categorically excludes from environmental review the action of
issuing an investigational exemption for a new drug intended to be administered to
human subjects, provided that the conditions of 21 CFR \S 25.15(c) and \S
25.15(d) are met and provided that no extraordinary circumstances exist that would
result in a significant environmental effect. The investigation described in this
IND meets these conditions. Daraxonrasib will be administered orally to a small
investigational population of up to 18 treated participants at a single academic
site, at clinical doses of 160 to 300 mg once daily over a limited perioperative
schedule; the quantities of drug substance and drug product manufactured,
distributed, and used are correspondingly small and are confined to the
controlled clinical setting and its regulated pharmaceutical-waste stream. The
estimated concentration of the drug substance entering the environment at the
expected use level is below the threshold at which an environmental assessment is
ordinarily required, and no extraordinary circumstances are present: the action
does not significantly affect the quality of the human environment, does not
involve a substance that is persistent, bioaccumulative, or otherwise of
recognized environmental concern at the quantities used, and does not affect any
species or habitat protected under federal environmental law
\cite{phase1guidance}.

The scale of material involved is bounded and can be stated explicitly. With a
maximum of eighteen treated participants distributed across the three dose levels
of the 3+3 design and a once-daily perioperative schedule, the total quantity of
daraxonrasib that could be administered over the entire investigation is on the
order of grams, not kilograms, and is dispensed, consumed, returned, or destroyed
entirely within the controlled clinical setting and its regulated
pharmaceutical-waste stream. A bounding arithmetic illustration confirms the
order of magnitude: even on the conservative assumption that every one of the
eighteen treated participants received the highest dose level of 300 mg once daily
for a continuous month of thirty days, the aggregate administered drug substance
would be eighteen times 300 mg times thirty days, equal to 162,000 mg, or about
0.16 kg, distributed across the entire study and across the perioperative schedule
in which dosing is paused around surgery rather than continuous. The realized
total is smaller still, because the lower dose levels deliver 160 mg and 220 mg
and because the schedule includes the pre-surgical pause and the restart
determined by the human-confirmed advisory rather than uninterrupted daily dosing.

This bounding quantity, entering the environment only as the small excreted
fraction after metabolism in a population confined to a single academic site and
managed through a regulated pharmaceutical-waste stream, is far below any
threshold at which an environmental assessment is ordinarily required, and it
confirms that the predicted environmental introduction concentration of the drug
substance is negligible. Investigational product that is unused, returned,
expired, or otherwise unsuitable is destroyed only after sponsor authorization
and in accordance with applicable federal, state, and local regulations governing
the disposal of pharmaceutical waste, and the destruction is documented in the
drug-accountability record, so that no meaningful quantity of drug substance
is released to the environment. 

The durable robotic device and its
single-use accessories are likewise managed within the regulated clinical-waste
stream and introduce no drug substance into the environment. To the sponsor's
knowledge, no data or information exist that would indicate that the proposed
investigation may significantly affect the quality of the human environment.
Accordingly, neither an environmental assessment nor an environmental impact
statement is required for this submission, and the categorical exclusion under 21
CFR \S 25.31(e) is claimed. Should the scope of the investigation expand in a
manner that would exceed the conditions of the categorical exclusion, for example
a substantial increase in the number of participants, the manufactured quantity,
or the duration of exposure, an environmental assessment under 21 CFR \S 25.40
would be submitted by information amendment before implementing the change.
